\section{Заключение}

В организационно-экономической части работы была выполнена оценка
трудоёмкости, стоимости разработки, затрат на внедрение, рыночных
параметров продукта и показателей инвестиционной эффективности
разрабатываемой методики верификации соответствия TLA+ спецификации и
исходного кода.

В результате проведённых расчётов установлено, что суммарная
трудоёмкость разработки составляет 436 чел.-ч, что соответствует
54.5 чел.-дням. Календарная продолжительность проекта по диаграмме
Ганта составляет 43 рабочих дня. Наиболее трудоёмким этапом является
разработка библиотеки инструментирования исходного кода, трудоёмкость
которой составляет 112 чел.-ч.

Расчёт затрат показал, что стоимость изготовления программного
продукта составляет 1443246.44 руб., а затраты на внедрение продукта
--- 2539632.90 руб. Следовательно, общие затраты на реализацию проекта
составляют 3982879.34 руб. Наибольшую долю в структуре затрат занимают
расходы на заработную плату исполнителей, оборудование и накладные
расходы.

В ходе исследования рынка установлено, что продукт ориентирован на
узкий сегмент корпоративных заказчиков, использующих формальные методы
верификации. Планируемый объём продаж составляет 15 копий в год, или
30 копий за двухлетний период. Принятая рыночная цена одной копии
продукта составляет 250000 руб. При этом годовая выручка прогнозируется
на уровне 3750000 руб., а выручка за два года --- 7500000 руб.

Расчёты показали, что полная себестоимость одной копии продукта
составляет 144789.70 руб., а прибыль от реализации одной копии ---
145200 руб. Годовая прибыль оценивается в 2178000 руб., а прибыль за
два года --- в 4356000 руб. Согласно балансу проекта, точка
безубыточности достигается примерно на 13--15 месяце эксплуатации.

Дополнительно была выполнена цифровая оценка инвестиций в инновационное
производство. При объёме инвестиций в изготовление продукта
$K = 1443246.44$ руб. получены следующие показатели: чистая приведённая
стоимость $NPV = 1693073.56$ руб., внутренняя норма доходности
$IRR = 119.62\%$, рентабельность инвестиций $ROI = 201.82\%$,
экономическая добавленная стоимость $EVA = 1889350.71$ руб., срок
окупаемости --- 7.95 месяцев.

Основные результаты организационно-экономического обоснования приведены
в таблице~\ref{tab:final_results}.

\begin{table}[h]
\centering
\caption{Итоговые технико-экономические показатели проекта}
\label{tab:final_results}
\resizebox{\textwidth}{!}{
\begin{tabular}{|p{10cm}|c|}
\hline
Показатель & Значение \\ \hline
Суммарная трудоёмкость разработки & 436 чел.-ч \\ \hline
Трудоёмкость разработки в человеко-днях & 54.5 чел.-дня \\ \hline
Календарная продолжительность проекта & 43 рабочих дня \\ \hline
Наиболее трудоёмкий этап & Разработка библиотеки трассировки, 112 чел.-ч \\ \hline
Затраты на изготовление продукта & 1443246.44 руб. \\ \hline
Затраты на внедрение продукта & 2539632.90 руб. \\ \hline
Общие затраты на реализацию проекта & 3982879.34 руб. \\ \hline
Цена одной копии продукта & 250000 руб. \\ \hline
Полная себестоимость одной копии & 144789.70 руб. \\ \hline
Планируемый объём продаж & 15 копий в год \\ \hline
Объём продаж за 2 года & 30 копий \\ \hline
Годовая выручка & 3750000 руб. \\ \hline
Выручка за 2 года & 7500000 руб. \\ \hline
Прибыль от одной копии & 145200 руб. \\ \hline
Годовая прибыль & 2178000 руб. \\ \hline
Прибыль за 2 года & 4356000 руб. \\ \hline
Срок достижения точки безубыточности по балансу & 13--15 месяцев \\ \hline
Чистая приведённая стоимость $NPV$ & 1693073.56 руб. \\ \hline
Внутренняя норма доходности $IRR$ & 119.62\% \\ \hline
Рентабельность инвестиций $ROI$ & 201.82\% \\ \hline
Экономическая добавленная стоимость $EVA$ & 1889350.71 руб. \\ \hline
Срок окупаемости инвестиций & 7.95 месяцев \\ \hline
\end{tabular}
}
\end{table}

Таким образом, выполненное организационно-экономическое обоснование
показывает, что разработка методики и программного инструмента
верификации является экономически целесообразной. Несмотря на
ограниченный рынок сбыта, специализированный характер продукта
обеспечивает достаточный уровень цены, позволяющий покрыть затраты на
разработку и внедрение и получить положительный финансовый результат.
